\section{\texorpdfstring{\(\ODBS\)}{O\_D\^BS} Register Blocks}

The paper contract is
\[
  \ODBS |0\rangle^{n-\ell}|s\rangle^\ell |i\rangle^n
  =
  |r_{si}\rangle^n |i\rangle^n .
\]
QBE stores the row register in bits \([1,1+n)\) and the full
\(\ODBS\) address register in bits \([1+n,1+2n)\).

\subsection{Executable Image}

For a compound basis index \(j\), define
\[
  \operatorname{row}(j)=(j\gg1)\mathbin{\&}(2^n-1),
\]
\[
  \operatorname{od}(j)=(j\gg(1+n))\mathbin{\&}(2^n-1),
\]
and
\[
  r(j)=\operatorname{robinSparseColumnMap}
  (n,\operatorname{sparse}(j),\operatorname{row}(j)).
\]
The executable paper-image skeleton is
\[
  F(j)=j\bmod 2^{1+n}
    + r(j)2^{1+n}
    + \left\lfloor\frac{j}{2^{1+2n}}\right\rfloor 2^{1+2n}.
\]

\subsection{Address Range}

Lean proves that \(\operatorname{row}(j)<2^n\) as
\Lean{bandedSparseAccessPaperRegisters\_row\_lt\_gridSize}.  Under the explicit
condition \(2\le n\), Lean proves
\[
  r(j)<2^n.
\]
This is \Lean{bandedSparseAccessPaperAddress\_lt\_gridSize\_of\_two\_le}; the
Boolean version is
\Lean{bandedSparseAccessPaperAddressInRange\_eq\_true\_of\_two\_le}.

\subsection{Image Range and Roundtrip}

Assume \(j\) is in the full finite basis and \(r(j)<2^n\).  The low splice is
below the next high-tail boundary:
\[
  j\bmod 2^{1+n}+r(j)2^{1+n}<2^{1+2n}.
\]
It follows that \(F(j)\) is still in the full finite basis.  Lean records this
as
\[
  \Lean{bandedSparseAccessPaperImage\_lt\_qubitDim\_of\_address\_lt}.
\]

The low block is preserved:
\[
  F(j)\bmod 2^{1+n}=j\bmod 2^{1+n}.
\]
Therefore
\[
  \operatorname{row}(F(j))=\operatorname{row}(j),
\]
which is \Lean{bandedSparseAccessPaperImage\_rowValue\_eq}.  Under the same
\(r(j)<2^n\) hypothesis,
\[
  \operatorname{od}(F(j))=r(j),
\]
which is \Lean{bandedSparseAccessPaperImage\_odRegisterValue\_eq}.

\subsection{No-Spill}

Let
\[
  H(j)=j\gg(1+2n).
\]
If \(r(j)<2^n\), the image splice does not carry into the high tail:
\[
  H(F(j))=H(j).
\]
This is \Lean{bandedSparseAccessPaperImage\_highTail\_eq\_of\_address\_lt}.
The Boolean no-spill theorem is
\Lean{bandedSparseAccessPaperImageNoSpill\_eq\_true\_of\_address\_lt}.

\subsection{Matrix Entry Bridge}

The forward skeleton is
\[
  M_{i,j}=
  \begin{cases}
  1, & i=F(j),\\
  0, & i\ne F(j).
  \end{cases}
\]
After packaging \(F(j)\) as a finite index, Lean proves
\[
  M_{F(j),j}=1
\]
as \Lean{bandedSparseAccessPaperMatrix\_imageFin\_eq\_one}.  The active gate
version is \Lean{oneTermRobinGate\_O\_D\_BS\_imageFin\_eq\_one}.  The dagger
skeleton satisfies
\[
  M^\dagger_{j,F(j)}=1,
\]
proved by \Lean{bandedSparseAccessPaperDaggerMatrix\_imageFin\_eq\_one}.

These are entry and register-safety blocks.  They do not prove injectivity,
inverse-on-range, cleanup, or unitarity.
